% =====================================================================
% Lightweight Protocol-Translation Bridges for Heterogeneous LLM
% Tool-Calling APIs: A Case Study on Codex-Ollama Interoperation
%
% arXiv-ready LaTeX source
% =====================================================================

\documentclass[11pt,a4paper]{article}

% --- Packages ---
\usepackage[utf8]{inputenc}
\usepackage[T1]{fontenc}
\usepackage[margin=1in]{geometry}
\usepackage{hyperref}
\usepackage{graphicx}
\usepackage{listings}
\usepackage{xcolor}
\usepackage{booktabs}
\usepackage{amsmath}
\usepackage{amssymb}
\usepackage{enumitem}
\usepackage{titling}
\usepackage{abstract}
\usepackage{float}

% --- Metadata ---
\title{\bf Lightweight Protocol-Translation Bridges for\\
Heterogeneous LLM Tool-Calling APIs:\\
A Case Study on Codex-Ollama Interoperation}

\author{
  xuanyuan\\
  \texttt{xrz5785@gmail.com}
}

\date{May 2026\\Technical Report v1.0}

% --- Code listing style ---
\definecolor{codebg}{rgb}{0.95,0.95,0.95}
\definecolor{codeframe}{rgb}{0.8,0.8,0.8}
\lstset{
  basicstyle=\ttfamily\small,
  backgroundcolor=\color{codebg},
  frame=single,
  rulecolor=\color{codeframe},
  breaklines=true,
  postbreak=\mbox{\textcolor{red}{$\hookrightarrow$}\space},
  showstringspaces=false,
  columns=flexible,
  aboveskip=8pt,
  belowskip=8pt,
}

% --- Compact itemize ---
\setlist{nosep,leftmargin=*}

\begin{document}

\maketitle

\begin{abstract}
Large Language Model (LLM) inference frameworks widely claim ``OpenAI API
compatibility,'' yet this compatibility often exists only at the syntactic
level---HTTP endpoints accept request payloads and return 200, but fail to
preserve semantic contracts around tool-calling (function-calling) behavior.
We document a concrete failure case: Ollama's \texttt{/v1/responses}
endpoint, when used as a provider for Codex CLI (OpenAI's coding agent),
returns malformed responses that cause all tool-call attempts to fail with
``unsupported call'' errors, despite the same models producing correct
\texttt{tool\_calls} through Ollama's \texttt{/v1/chat/completions}
endpoint.

We propose the \textbf{Protocol-Translation Bridge (PTB)} pattern: a
lightweight, zero-configuration proxy that performs three core
transformations---(1) request format translation from Responses API to Chat
Completions API, (2) tool schema simplification to fit local models'
attention budgets, and (3) reverse synthesis of SSE event streams from
non-streaming upstream responses. Our implementation ($\sim$800 lines of
Python) restores full tool-calling functionality across all tested local
models (qwen3:14b, huihui4:8b-a4b, qwen2.5-coder:3b) with zero
modifications to either the client framework or the inference engine.

The key insight is counterintuitive: \textit{downgrading} from the newer
Responses API to the older Chat Completions API increases semantic fidelity,
because the older endpoint has more mature native implementations within
inference frameworks.
\end{abstract}

% =====================================================================
\section{Introduction}
% =====================================================================

\subsection{Background}

LLM-based coding agents---Codex \cite{codex}, Cursor, Copilot,
Aider---rely on tool-calling APIs to execute shell commands, read and write
files, spawn sub-agents, and interact with users. These agents typically
target OpenAI's API contract, either the Chat Completions API
(\texttt{/v1/chat/completions}) or the newer Responses API
(\texttt{/v1/responses}) \cite{openai_responses}. The Responses API is the
recommended modern endpoint for agent frameworks because it provides
structured multi-output responses, native streaming SSE events, and a
unified interface for text and tool-call outputs.

Meanwhile, the local-model ecosystem---Ollama \cite{ollama}, vLLM
\cite{vllm}, llama.cpp---has rapidly evolved to support tool-calling.
Ollama's documentation states that its \texttt{/v1/responses} endpoint is
``OpenAI-compatible'' \cite{ollama_compat}. Codex CLI v0.130.0 added
\texttt{--oss --local-provider ollama} flags to support local models through
this endpoint.

\subsection{The Problem}

When Codex connects to Ollama's \texttt{/v1/responses} endpoint with a
local model, tool calls fail with:

\begin{lstlisting}
unsupported call: call_<id> (exec_command)
\end{lstlisting}

The error is silent at the HTTP level---Ollama returns 200 OK. The failure
occurs at the semantic level: the response body does not contain the
structured \texttt{function\_call} output items that Codex's agent runtime
expects. The same model, queried via \texttt{/v1/chat/completions} with
identical tool definitions, produces correct \texttt{tool\_calls} in its
response.

This paper documents the diagnosis, solution, and generalization of this
problem.

% =====================================================================
\section{Problem Analysis}
% =====================================================================

\subsection{Root Cause}

Ollama's \texttt{/v1/responses} endpoint is a thin wrapper around its native
\texttt{/api/chat} endpoint. The internal path is:

\begin{lstlisting}
Codex -> POST /v1/responses (Ollama's OpenAI-compat layer)
      -> Ollama rewraps as /api/chat request
      -> Model outputs text (not structured function_call)
      -> Ollama returns {output: [{type: "message", content: "..."}]}
      -> Codex finds no function_call -> "unsupported call"
\end{lstlisting}

When the same model is queried via \texttt{/v1/chat/completions}:

\begin{lstlisting}
Client -> POST /v1/chat/completions
       -> Ollama rewraps as /api/chat with tools properly nested
       -> Model outputs {tool_calls: [{function: {name, arguments}}]}
       -> Ollama returns {choices: [{message: {tool_calls: [...]}}]}
\end{lstlisting}

The critical difference: Ollama's native \texttt{/api/chat} endpoint
natively supports tool-calling and correctly passes tool definitions to the
model. The \texttt{/v1/responses} wrapper loses this capability during
format translation.

\subsection{Why Not Fix Upstream?}

\begin{enumerate}[nosep]
  \item Ollama's issue tracker already has reports about
        \texttt{/v1/responses} tool-calling gaps
  \item Fix timeline unknown---the endpoint is not Ollama's priority
  \item A proxy is zero-friction---no need to modify Codex or Ollama
  \item Generalizes---the same pattern applies to any pair of incompatible
        LLM API surfaces
\end{enumerate}

\subsection{Diagnosis Methodology}

We employed a 12-step iterative diagnosis, summarized in Table~\ref{tab:diagnosis}.

\begin{table}[H]
\centering
\caption{12-Step Iterative Diagnosis}
\label{tab:diagnosis}
\begin{tabular}{ccll}
\toprule
\textbf{Step} & \textbf{Hypothesis} & \textbf{Test} & \textbf{Result} \\
\midrule
1  & Direct \texttt{/v1/responses} works    & Codex + Ollama                  & ``unsupported call'' \\
2  & Model capability issue                 & Manual \texttt{/v1/chat/completions} & Tool calls work \\
3  & Ollama responses impl broken           & Read source                     & Thin wrapper, loses tools \\
4  & Tool choice mode                       & \texttt{required} vs \texttt{auto} & \texttt{required} causes loops \\
5  & Too many tools                         & 11 tools $\rightarrow$ 3 tools  & Better, still inconsistent \\
6  & Tool parameter overload                & 10 params $\rightarrow$ 2--3  & Model selects correctly \\
7  & Model outputs JSON text                & Check response format           & Text JSON, not \texttt{tool\_calls} \\
8  & System prompt too weak                 & Add CRITICAL directive          & Model starts calling tools \\
9  & Usage field mismatch                   & Check \texttt{input\_tokens}    & Missing, causes disconnect \\
10 & \texttt{output\_index} hardcoded       & Multi-output responses          & Text overwrites \texttt{function\_call} \\
11 & Pull fails for custom models           & POST \texttt{/api/pull}         & ``file does not exist'' \\
12 & GGUF architecture compatibility        & gemma4 in Ollama                & Needs specific arch support \\
\bottomrule
\end{tabular}
\end{table}

% =====================================================================
\section{The Protocol-Translation Bridge (PTB) Pattern}
% =====================================================================

\subsection{Architecture}

The proxy sits between Codex and Ollama, listening on port 11434 with
Ollama on 11433:

\begin{figure}[H]
\centering
\fbox{
\begin{minipage}{0.92\textwidth}
\centering
\vspace{3pt}
\textbf{Figure 1: PTB Architecture}\\[8pt]
\begin{tabular}{c c c c c}
\texttt{Codex} & $\xrightarrow[\text{SSE}]{\text{/v1/responses}}$ &
\texttt{Bridge :11434} & $\xrightarrow[\text{JSON}]{\text{/v1/chat/completions}}$ &
\texttt{Ollama :11433} \\
& $\xleftarrow[\text{events}]{\text{}}$ &
& $\xleftarrow[\text{response}]{\text{}}$ &
\end{tabular}
\vspace{6pt}
\end{minipage}
}
\end{figure}

The proxy is transparent for non-Responses requests (pass-through for
\texttt{/api/tags}, \texttt{/api/chat}, etc.), and only intervenes for two
specific paths: \texttt{POST /v1/responses} (full protocol translation) and
\texttt{POST /api/pull} (intercept for locally-available models).

\subsection{Transformation 1: Request Format Translation}

The Responses API request is converted to Chat Completions format.
Key design decisions:

\begin{itemize}
  \item \texttt{stream: false}---we always request non-streaming from
        Ollama, then synthesize SSE events ourselves. This avoids the
        complexity of real-time SSE$\rightarrow$SSE translation and gives us
        complete control over event ordering.
  \item \texttt{tool\_choice: "auto"}---\texttt{"required"} causes infinite
        tool-call loops with some models. \texttt{"auto"} combined with a
        strong system prompt provides the right balance.
  \item \texttt{instructions} becomes part of the system message, augmented
        with tool-usage directives.
\end{itemize}

\subsection{Transformation 2: Tool Schema Simplification}

Codex's internal tools are complex: \texttt{exec\_command} alone has 10
parameters. Across 11 tools, the total tool definition is approximately
4,100 tokens---exceeding the effective attention budget of 8B-class models.

\begin{table}[H]
\centering
\caption{Tool Parameter Reduction}
\label{tab:tools}
\begin{tabular}{lcc}
\toprule
\textbf{Tool} & \textbf{Original Params} & \textbf{Essential Params} \\
\midrule
\texttt{exec\_command}    & 10 & 2 (cmd, workdir) \\
\texttt{write\_stdin}     & 6  & 2 (session\_id, chars) \\
\texttt{spawn\_agent}     & 8  & 3 (agent\_type, items, message) \\
\texttt{view\_image}      & 3  & 1 (path) \\
\texttt{update\_plan}     & 2  & 1 (plan) \\
\texttt{request\_user\_input} & 3 & 1 (questions) \\
\texttt{send\_input}      & 4  & 3 (target, message, items) \\
\texttt{resume\_agent}    & 2  & 1 (id) \\
\texttt{wait\_agent}      & 3  & 1 (targets) \\
\texttt{close\_agent}     & 2  & 1 (target) \\
\bottomrule
\end{tabular}
\end{table}

After simplification: $\sim$800 tokens (5$\times$ reduction). Codex's tool
executor fills in sensible defaults for omitted parameters. This is a
\textit{zero-cost accuracy improvement}---the model selects the correct tool
with higher probability because the signal-to-noise ratio in the tool
definitions is higher.

\subsection{Transformation 3: Response Event Synthesis}

From a non-streaming Chat Completions JSON response, we synthesize the
OpenAI Responses API SSE event stream:

\begin{figure}[H]
\centering
\fbox{
\begin{minipage}{0.9\textwidth}
\vspace{3em}
\centering
\textbf{Figure 2: SSE Event Synthesis Flow}\\[6pt]
\texttt{chat.completion JSON} $\rightarrow$
\texttt{response.created} $\rightarrow$
\texttt{in\_progress} $\rightarrow$\\[2pt]
\texttt{output\_item.added (function\_call)} $\rightarrow$
\texttt{function\_call\_arguments.delta} $\rightarrow$\\
\texttt{function\_call\_arguments.done} $\rightarrow$
\texttt{output\_item.done} $\rightarrow$\\
\texttt{output\_item.added (message)} $\rightarrow$
\texttt{content\_part.added} $\rightarrow$\\
\texttt{output\_text.delta (xN)} $\rightarrow$
\texttt{output\_item.done} $\rightarrow$\\
\texttt{response.completed (with normalized usage)}
\vspace{3em}
\end{minipage}
}
\end{figure}

Critical details that caused failures:

\begin{enumerate}[nosep]
  \item \texttt{output\_index}: MUST be 0 for the first function\_call, 1
        for the text message. Hardcoding 0 causes multi-output corruption.
  \item \texttt{sequence\_number}: MUST be globally monotonically increasing
        across all events.
  \item \texttt{usage}: Ollama returns \texttt{\{prompt\_tokens,
        completion\_tokens\}} but Codex expects \texttt{\{input\_tokens,
        output\_tokens\}}.
  \item Event ordering is a strict protocol; deviations cause client
        disconnections.
\end{enumerate}

\subsection{Pull Interception}

Codex calls \texttt{POST /api/pull} for every model before first use.
Custom GGUF models (e.g., huihui4-8b-a4b) are not in Ollama's registry,
causing pull failures. The proxy intercepts this endpoint, checks if the
model exists locally via \texttt{/api/tags}, and returns
\texttt{\{"status":"success"\}\textbackslash n} (NDJSON format) for
locally-available models.

% =====================================================================
\section{Implementation}
% =====================================================================

\subsection{Technology Stack}

\begin{itemize}[nosep]
  \item \textbf{Language}: Python 3.12 (standard library + aiohttp)
  \item \textbf{Lines of code}: $\sim$480 effective, 807 total
  \item \textbf{Dependencies}: aiohttp only
  \item \textbf{Deployment}: macOS launchd (KeepAlive daemon) or manual
        \texttt{python3 proxy.py}
\end{itemize}

\subsection{Code Organization}

\begin{lstlisting}
proxy.py
|-- normalize_usage()          -- Usage field mapping (Ollama -> OpenAI)
|-- simplify_tools()           -- Tool parameter reduction
|-- responses_to_chat()        -- Request format translation
|-- SSEResponseBuilder         -- SSE event synthesis engine
|   |-- start() / in_progress() / complete() / error()
|   |-- add_text_delta()
|   |-- add_tool_call_start() / add_tool_args_delta() / finish_tool_call()
|-- proxy_handler()            -- Main request dispatcher
|-- _synthesize_sse()          -- SSE stream construction
|-- health_handler()           -- Health check endpoint
|-- main()                     -- CLI, signal handling, startup check
\end{lstlisting}

\subsection{Deployment}

Production deployment uses macOS launchd with KeepAlive:

\begin{lstlisting}[language=XML]
<!-- ~/Library/LaunchAgents/com.x.codex-bridge.plist -->
<key>RunAtLoad</key><true/>
<key>KeepAlive</key><true/>
\end{lstlisting}

A control script provides standard service management:

\begin{lstlisting}[language=bash]
codex-bridge-ctl.sh start|stop|restart|status|logs
\end{lstlisting}

Codex is configured with shell aliases:

\begin{lstlisting}[language=bash]
alias cx14='codex --oss --local-provider ollama -m qwen3:14b'
alias cx14e='codex exec --skip-git-repo-check --oss
             --local-provider ollama -m qwen3:14b'
\end{lstlisting}

% =====================================================================
\section{Experimental Validation}
% =====================================================================

\subsection{Test Setup}

\begin{table}[H]
\centering
\caption{Test Environment}
\label{tab:env}
\begin{tabular}{ll}
\toprule
\textbf{Component} & \textbf{Version} \\
\midrule
Codex CLI  & v0.130.0 / v0.131.0 \\
Ollama     & 0.23.4 \\
Bridge     & v1.1.0 \\
OS         & macOS 26.4 (Apple Silicon) \\
\bottomrule
\end{tabular}
\end{table}

\subsection{Model Compatibility}

\begin{table}[H]
\centering
\caption{Model Compatibility Matrix}
\label{tab:models}
\begin{tabular}{lcccc}
\toprule
\textbf{Model} & \textbf{Size} & \textbf{Tools} & \textbf{Chinese} & \textbf{Notes} \\
\midrule
qwen3:14b              & 9.3GB  & Stable  & Native & Flagship \\
huihui4:8b-a4b         & 5.4GB  & Good    & Yes    & MoE, 4/8.1B active \\
Qwen2.5-Coder-7B-GGUF  & 7B     & Moderate& Yes    & Backup \\
qwen2.5-coder:3b       & 1.9GB  & Weak    & Yes    & Lightweight text \\
llama3.1:8b            & 4.9GB  & Weak    & No     & English only \\
\bottomrule
\end{tabular}
\end{table}

\subsection{End-to-End Test}

\begin{lstlisting}
$ codex exec --skip-git-repo-check --ephemeral --oss \
    --local-provider ollama -m "huihui4-8b-a4b:latest" \
    "list files in proxy/ directory"

exec
/bin/zsh -lc 'ls -R /Users/x/ai-assets/codex-proxy/'
  succeeded in 0ms:
__pycache__
proxy.py
paper/

/Users/x/ai-assets/codex-proxy/__pycache__:
proxy.cpython-312.pyc
\end{lstlisting}

\subsection{Failure Analysis}

During development, we encountered and resolved 10 distinct failure modes,
summarized in Table~\ref{tab:failures}.

\begin{table}[H]
\centering
\caption{Failure Mode Analysis (10 Error Types)}
\label{tab:failures}
\begin{tabular}{cll}
\toprule
\textbf{\#} & \textbf{Failure} & \textbf{Root Cause \& Fix} \\
\midrule
1  & Port conflict                     & Ollama launchd auto-restart; manual port mgmt \\
2  & HTTP 400 Bad Request              & Content-Length not updated after body change \\
3  & Content-Length off by 1           & stream:false changed after header calc; reorder \\
4  & Empty function\_call name         & 4100-token tool def overload; simplify\_tools() \\
5  & Text JSON instead of tool\_calls  & Model outputs JSON as text; enhanced system prompt \\
6  & Missing input\_tokens             & Usage field naming mismatch; normalize\_usage() \\
7  & Transport closed                  & Malformed response.completed; fix usage \\
8  & output\_index hardcoded to 0      & Multi-output ordering; dynamic output\_index \\
9  & Pull failure for custom models    & Not in registry; pull interception \\
10 & Literal \textbackslash n in pull  & Escaped vs real newline; binary-correct \textbackslash n \\
\bottomrule
\end{tabular}
\end{table}

% =====================================================================
\section{Discussion}
% =====================================================================

\subsection{The ``Newer API is Better'' Fallacy}

A counterintuitive finding: the newer \texttt{/v1/responses} endpoint
(introduced by OpenAI in 2025) performed \textit{worse} than the older
\texttt{/v1/chat/completions} endpoint for tool-calling through Ollama.
This is because the Chat Completions API has been the primary integration
target for inference frameworks for years, receiving more testing and native
optimization. The Responses API, being newer, has thinner compatibility
wrappers.

\textbf{Lesson:} When debugging API compatibility issues, try downgrading to
an older API surface before assuming the model or framework is broken.

\subsection{Attention Budget as a First-Class Constraint}

Tool definitions consume prompt tokens. For an 8B model with a 32K context
window, 4,100 tokens of tool definitions represent $\sim$13\% of total
budget. But the effective attention budget for tool selection is much
smaller---the model must attend to the system prompt, conversation history,
\textit{and} tool definitions simultaneously.

Our simplification from 4,100 $\rightarrow$ 800 tokens (5$\times$ reduction)
was the single most impactful change for model accuracy. This suggests that
tool definition design for local models should be treated as a prompt
engineering problem, not just an API integration problem.

\subsection{SSE Synthesis vs. Real-Time Translation}

\begin{table}[H]
\centering
\caption{Trade-off: SSE Synthesis vs. Real-Time Translation}
\begin{tabular}{lcc}
\toprule
& \textbf{Real-Time SSE} & \textbf{Synthesized SSE} \\
\midrule
Latency       & Low (streaming)  & First-byte $\approx$ generation time \\
Complexity    & High (state machine) & Low ($\sim$480 loc) \\
Correctness   & Event reordering risk & Guaranteed correct ordering \\
Best for      & Fast models, production & Local models, prototyping \\
\bottomrule
\end{tabular}
\end{table}

For local models where generation latency is typically 5--30 seconds, the
first-byte latency of non-streaming is acceptable. For production deployments
with faster models, real-time translation would be the next optimization.

\subsection{Generalizability}

The PTB pattern applies beyond Codex-Ollama. Any pair of LLM API surfaces
with syntactic-but-not-semantic compatibility can be bridged: Cursor +
Ollama, Continue.dev + vLLM, LangChain agents + llama.cpp. The core
principle: \textbf{identify the API surface where tool-calling works
natively, translate requests to that surface, and translate responses back
to the client's expected surface.}

\subsection{Limitations}

\begin{enumerate}[nosep]
  \item No real streaming---first-byte latency equals full generation time
  \item Single-model focus---no load balancing across multiple instances
  \item No authentication---assumes local-only deployment
  \item No automated tests---manual end-to-end testing only
  \item No response caching---repeated queries regenerate
\end{enumerate}

% =====================================================================
\section{Related Work}
% =====================================================================

\begin{itemize}
  \item \textbf{LiteLLM} \cite{litellm}: Universal LLM proxy supporting
        $\sim$100 providers. Primarily targets Chat Completions; Responses
        API support is nascent.
  \item \textbf{OpenRouter}: Commercial routing service. Does not address
        Responses API.
  \item \textbf{vLLM} \cite{vllm}: OpenAI-compatible server. Focused on
        serving, not protocol translation between API surfaces.
  \item \textbf{OpenAI Agents SDK}: Official agent framework using
        Responses API. Our work enables these agents with local models.
  \item \textbf{porter} \cite{porter}: Lightweight LLM proxy for
        authentication and routing, not protocol translation.
\end{itemize}

The unique contribution of this work is the combination of: (a) Responses
API $\leftrightarrow$ Chat Completions semantic translation, (b) tool
schema simplification for local models, and (c) reverse SSE synthesis from
non-streaming responses.

% =====================================================================
\section{Conclusion}
% =====================================================================

We presented the Protocol-Translation Bridge (PTB) pattern, a lightweight
solution to the problem of heterogeneous LLM API tool-calling compatibility.
Our implementation for Codex-Ollama interoperation ($\sim$800 lines of
Python) successfully restores tool-calling functionality for local models,
with zero modifications to either the client framework or the inference
engine.

Key findings:
\begin{enumerate}[nosep]
  \item API ``compatibility'' claims require semantic-level verification
  \item Downgrading to older API surfaces can increase semantic fidelity
  \item Tool schema simplification is a zero-cost optimization for local models
  \item SSE synthesis from non-streaming responses is a viable alternative
        to real-time SSE translation
  \item Usage field naming varies across implementations and requires normalization
\end{enumerate}

The code, deployment configuration, and this paper are released as
open-source at \texttt{/Users/x/ai-assets/codex-proxy/}.

% =====================================================================
\bibliographystyle{plain}
\begin{thebibliography}{99}

\bibitem{openai_responses}
OpenAI. ``Responses API Reference.''
\texttt{https://platform.openai.com/docs/api-reference/responses}

\bibitem{codex}
OpenAI. ``Codex CLI.''
\texttt{https://github.com/openai/codex}

\bibitem{ollama}
Ollama. ``Get up and running with large language models.''
\texttt{https://ollama.com}

\bibitem{ollama_compat}
Ollama. ``OpenAI Compatibility.''
\texttt{https://ollama.com/blog/openai-compatibility}

\bibitem{litellm}
BerriAI. ``LiteLLM: Call all LLM APIs using the OpenAI format.''
\texttt{https://github.com/BerriAI/litellm}

\bibitem{vllm}
vLLM Team. ``vLLM: Easy, Fast, and Cheap LLM Serving for Everyone.''
\texttt{https://docs.vllm.ai}

\bibitem{huihui4}
Huihui-ai. ``Huihui4-8B-A4B: Mixture-of-Experts Language Model.''
\texttt{https://huggingface.co/huihui-ai/Huihui4-8B-A4B-GGUF}

\bibitem{qwen3}
Qwen Team. ``Qwen3: Technical Report.'' 2025.

\bibitem{claude_code}
Anthropic. ``Claude Code: Agentic coding tool.''
\texttt{https://docs.anthropic.com/en/docs/claude-code}

\bibitem{porter}
Porter.sh. ``Lightweight LLM API Proxy.''
\texttt{https://porter.sh}

\end{thebibliography}

\end{document}
